\documentclass[a4paper]{article}

\usepackage{graphicx}
\usepackage{amsmath,amssymb}
\usepackage{minted}
\usepackage{multirow}
\usepackage{float}
\usepackage{algorithm}
\usepackage{algpseudocode}
\usepackage{todonotes}
\usepackage{forest}
\usepackage{xstring}
\usepackage{xcolor}
\usepackage[most]{tcolorbox}
\usepackage{xparse}
\usepackage{natbib}
\usepackage{adjustbox}

\usetikzlibrary{graphs,graphdrawing}
\usegdlibrary{layered}

% for external graphviz
\input{/home/donkarlo/Dropbox/repo/nd_language_ai_project/src/nd_language_ai/formal/computer/programming/tex/latex/code_clips/external_graphviz.tex}

% for tags
\input{/home/donkarlo/Dropbox/repo/nd_language_ai_project/src/nd_language_ai/formal/computer/programming/tex/latex/part/control_sequence/kind/semtag/semtag.tex}

% for links
\input{/home/donkarlo/Dropbox/repo/nd_language_ai_project/src/nd_language_ai/formal/computer/programming/tex/latex/code_clips/seehere.tex}

% for nested lists and sections
\input{/home/donkarlo/Dropbox/repo/nd_language_ai_project/src/nd_language_ai/formal/computer/programming/tex/latex/code_clips/deep_structure.tex}

\title{Controller}
\date{}
\author{Mohammad Rahmani}

\begin{document}
   \maketitle

   A controller is the part of a control system that determines the control input applied to a plant or process in order to influence its behavior toward a desired objective \cite{astrom-2008-feedback-systems-an-introduction-for-scientists-and-engineers}. Depending on the control architecture, the controller may use a reference or setpoint, measured outputs, estimated states, past errors, or a mathematical model of the plant.

   \section{Controller in a control system}
      Let $r(t)$ denote a desired reference, $y(t)$ the measured output of the plant, and $u(t)$ the control input generated by the controller.

      In a feedback controller, the control error is commonly defined as
      \[
         e(t)=r(t)-y(t).
      \]

      The controller then applies a control law that maps available information to the control input,
      \[
         u(t)=\mathcal{C}(r(t),y(t),e(t),\ldots),
      \]
      where $\mathcal{C}$ represents the controller. The exact form of this mapping depends on the controller type \cite{astrom-2008-feedback-systems-an-introduction-for-scientists-and-engineers}.

   \section{Open-loop and closed-loop control}
      In open-loop control, the controller generates the command without using feedback from the measured output. Therefore, the command is not corrected directly according to the actual response of the plant.

      In closed-loop or feedback control, measurements of the plant output are returned to the controller. The controller can then modify the control input according to the difference between desired and actual behavior \cite{astrom-2008-feedback-systems-an-introduction-for-scientists-and-engineers}.

      \seeherefooter{/home/donkarlo/Dropbox/repo/nd_robotic_ai_learning_project/src/control_theory/controller/open_loop_controller.tex}
      \seeherefooter{/home/donkarlo/Dropbox/repo/nd_robotic_ai_learning_project/src/control_theory/controller/closed_loop.tex}

   \section{Common controller types}
      \paragraph{Proportional controller}
      A proportional controller generates a control input proportional to the current error,
      \[
         u(t)=K_P e(t).
      \]
      It is a direct control law and does not require an online optimization problem \cite{astrom-2008-feedback-systems-an-introduction-for-scientists-and-engineers}.

      \seeherefooter{/home/donkarlo/Dropbox/repo/nd_robotic_ai_learning_project/src/control_theory/kind/proportional_controller/proportional_controller.tex}

      \paragraph{Proportional--integral--derivative controller}
      A PID controller combines proportional, integral, and derivative terms of the error,
      \[
         u(t)
         =
         K_P e(t)
         +
         K_I\int_0^t e(\tau)\,d\tau
         +
         K_D\frac{de(t)}{dt}.
      \]
      Like proportional control, a conventional PID controller applies a predefined control law rather than solving an online cost-minimization problem at each control step \cite{astrom-2008-feedback-systems-an-introduction-for-scientists-and-engineers}.

      \seeherefooter{/home/donkarlo/Dropbox/repo/nd_robotic_ai_learning_project/src/control_theory/kind/proportional_integral_derivative_controller/proportional_integral_derivative_controller.tex}

      \paragraph{Model predictive control}
      Model predictive control uses a model of the plant to predict future behavior and repeatedly solves a finite-horizon optimal-control problem. The first control action of the optimized sequence is applied, and the optimization is repeated at the next sampling step \cite{mayne-2000-constrained-model-predictive-control-stability-and-optimality}.

      \seeherefooter{/home/donkarlo/Dropbox/repo/nd_robotic_ai_learning_project/src/control_theory/kind/model_predictive_control/model_predictive_control.tex}

   \section{Cascade control}
      Controllers can also be connected in a cascade. In cascade control, the output of an outer controller becomes the setpoint of an inner controller \cite{lee-1998-pid-controller-tuning-for-cascade-control-systems}. For example, an outer proportional position controller can generate a desired velocity that is tracked by an inner PID velocity controller.

      \seeherefooter{/home/donkarlo/Dropbox/repo/nd_robotic_ai_learning_project/src/control_theory/controller/cascade_controller/cascade_controller.tex}

   \bibliographystyle{plainnat}
   \bibliography{/home/donkarlo/Dropbox/repo/nd_shared_research/bibliography/bibliography.bib}
\end{document}
